% Options for packages loaded elsewhere
\PassOptionsToPackage{unicode}{hyperref}
\PassOptionsToPackage{hyphens}{url}
%
\documentclass[
]{article}
\usepackage{amsmath,amssymb}
\usepackage{iftex}
\ifPDFTeX
  \usepackage[T1]{fontenc}
  \usepackage[utf8]{inputenc}
  \usepackage{textcomp} % provide euro and other symbols
\else % if luatex or xetex
  \usepackage{unicode-math} % this also loads fontspec
  \defaultfontfeatures{Scale=MatchLowercase}
  \defaultfontfeatures[\rmfamily]{Ligatures=TeX,Scale=1}
\fi
\usepackage{lmodern}
\ifPDFTeX\else
  % xetex/luatex font selection
\fi
% Use upquote if available, for straight quotes in verbatim environments
\IfFileExists{upquote.sty}{\usepackage{upquote}}{}
\IfFileExists{microtype.sty}{% use microtype if available
  \usepackage[]{microtype}
  \UseMicrotypeSet[protrusion]{basicmath} % disable protrusion for tt fonts
}{}
\makeatletter
\@ifundefined{KOMAClassName}{% if non-KOMA class
  \IfFileExists{parskip.sty}{%
    \usepackage{parskip}
  }{% else
    \setlength{\parindent}{0pt}
    \setlength{\parskip}{6pt plus 2pt minus 1pt}}
}{% if KOMA class
  \KOMAoptions{parskip=half}}
\makeatother
\usepackage{xcolor}
\usepackage[margin=1in]{geometry}
\usepackage{graphicx}
\makeatletter
\def\maxwidth{\ifdim\Gin@nat@width>\linewidth\linewidth\else\Gin@nat@width\fi}
\def\maxheight{\ifdim\Gin@nat@height>\textheight\textheight\else\Gin@nat@height\fi}
\makeatother
% Scale images if necessary, so that they will not overflow the page
% margins by default, and it is still possible to overwrite the defaults
% using explicit options in \includegraphics[width, height, ...]{}
\setkeys{Gin}{width=\maxwidth,height=\maxheight,keepaspectratio}
% Set default figure placement to htbp
\makeatletter
\def\fps@figure{htbp}
\makeatother
\setlength{\emergencystretch}{3em} % prevent overfull lines
\providecommand{\tightlist}{%
  \setlength{\itemsep}{0pt}\setlength{\parskip}{0pt}}
\setcounter{secnumdepth}{5}
\ifLuaTeX
  \usepackage{selnolig}  % disable illegal ligatures
\fi
\usepackage{bookmark}
\IfFileExists{xurl.sty}{\usepackage{xurl}}{} % add URL line breaks if available
\urlstyle{same}
\hypersetup{
  hidelinks,
  pdfcreator={LaTeX via pandoc}}

\author{}
\date{\vspace{-2.5em}}

\begin{document}

\vspace*{\fill}

\begin{center}
\Huge \textbf{Modeling Cetacean Distribution Patterns Using Spatial-Temporal GLMMs} \\[0.5cm]
\large By Joshua Charfauros, Justin Lang, Peter Xiong, Valerie De La Fuente, Ziqian Zhao, Samantha Su \\[0.2cm]
\normalsize Sponsor: Michaela Alksne, Scripps Institute of Oceanography.\\
\normalsize Mentor: Professor Baracaldo

\end{center}
\vspace*{\fill}
\thispagestyle{empty}
\newpage

\section{Introduction}\label{introduction}

Researching the presence of blue and fin whales and their reliance on
distinct call types for behaviors (such as foraging and reproduction) is
important for monitoring ecosystem health and guiding conservation
efforts in the California Current System. By analyzing their calls in
relation with environmental data, we can gain insight into their
behavioral patterns and habitat use across time.

Several foundational studies guided our research by providing us a solid
background understanding on our topic. For example, Becker et al.~(2022)
studied how different whale species use habitat across the California
Current and showed that it's important to use flexible models that can
account for changes in space and time. Campbell et al.~(2015) showed
that there were long-term trends in whale sightings using CalCOFI visual
survey data, while Oleson et al.~(2007) and Širović et al.~(2013)
demonstrated that specific whale call types correlate with behavioral
states, supporting our decision to model each call type separately.

Therefore, the goal of this project is to develop species distribution
models for blue and fin whale acoustic presence using environmental
drivers, and to visualize their occurrence across time. We also aim to
interpret call-type-specific patterns to better understand behavioral
ecology.

\textbf{Github:} All work will be documented in the following GitHub
repository:\\
\url{https://github.com/Capstone-24-25/capstone-scripps}

\section{Problem(s) of Interest}\label{problems-of-interest}

The primary problem we want to address is how environmental conditions
affect the acoustic presence of blue and fin whales in the Southern
California Bight. We will focus on modeling this presence across time
using call type data to better understand how different behaviors (such
as feeding vs.~mating) are distributed in relation to the environment.

A second core problem is identifying how specific whale call types (A,
B, D, 20 Hz, 40 Hz) reflect distinct behavioral states. By visualizing,
animating, and modeling each call type, we hope to develop inferences on
patterns related to foraging or reproductive behavior and explore how
these patterns change seasonally and geographically.

\section{Materials and Methods}\label{materials-and-methods}

\subsection{Datasets}\label{datasets}

\subsubsection{CASE-STSE}\label{case-stse}

The CASE-STSE dataset collects oceanographic data over one- to
three-month periods to generate detailed analyses and 30-day forecasts
of ocean conditions. It includes observations from various sources,
including underwater gliders, temperature probes, floating sensors, and
satellites. The dataset includes three-dimensional fields of
temperature, salinity, pressure, and velocity, making it a valuable
resource for understanding dynamic ocean processes. More information and
data products can be accessed at
\url{https://ecco.ucsd.edu/case_stse_results1.html}.

\subsubsection{CalCoFI}\label{calcofi}

The CalCoFI dataset contains data on whale detection across survey
stations, comprising of visual surveys, acoustic recordings, and
environmental DNA detections. The resulting data includes the number of
individual calls and their types (or response variables), such as A, B,
and D calls for blue whales, and A, B, D, 20 Hz and 40 Hz calls for fin
whales. These response variables will tell us how environmental
variables affects different whales/each individual call. Prior to
conducting our project, we conducted exploratory data analysis to better
understand this dataset.

\subsection{Data Wrangling Processing
Techniques}\label{data-wrangling-processing-techniques}

Before modeling, we combined acoustic data from CalCOFI with
environmental data from CASE-STSE to relate whale acoustic presence to
environmental drivers across space and time in the Southern California
Bight. We adapted a provided MATLAB script to extract relevant CASE-STSE
variables, then computed cruise averages for each day with acoustic
data, resulting in a dataset indexed by year and month.

\textbf{Merging Procedure:}

\begin{itemize}
\item
  Temporal matching: We joined CalCOFI and CASE-STSE records using the
  common year-month field.
\item
  Spatial matching: For each CalCOFI record, we identified the closest
  CASE-STSE grid cell using K-nearest neighbors on lattitude and
  longitude coordinates.
\end{itemize}

\subsection{Animations and
Visualizations}\label{animations-and-visualizations}

Michaela Alksne provided us with a MATLAB script that would plot our
extracted environmental data and show the difference over time. We
adapted this script in order to show the call types over these data in a
form of animated visualization, giving us a better understanding of our
data prior to modeling.

To do this, we graphed the call types over environmental data at
specific depths. At each specific depth and date, we parsed through the
matching CSV using a custom function which converted the csv into
numerical matrices. Then, as a base layer for water velocity, used
MATLAB's quiver function to properly visualize how water flows
throughout certain latitudes.

Then, we plotted different colored dots for each call type over these
environmental data and increased the size of the dots based on how long
the calls were. This allowed us to visualize how different call types
appear across space and time, and how they relate to environmental
conditions. We also set a predetermined color for each call so that we
can show them on a legend to the right.

Finally, we set all these collapsed images into one singular frame and
binded all of them together through rgb2ind to create our finished gif
representing the whale calls overlayed on different environmental data
at different depths.

\subsection{Statistical Modeling}\label{statistical-modeling}

We modeled fin and blue whale acoustic presence using two statistical
frameworks: a \textbf{tweedie generalized linear model (GLM)}
(Appropriate for continuous data with a large number of zeros
(zero-inflated)) and a \textbf{delta-gamma hurdle model} (Separately
models the probability of presence and the magnitude of response when
present), both implemented using predictor (written by Jim Thorson, a
collaborator on this project). The response variable, fin\_resp, was
constructed as the sum of 20 Hz and 40 Hz fin whale call rates,
representing total acoustic activity.

\subsubsection{Predictor Selection}\label{predictor-selection}

\begin{itemize}
\tightlist
\item
  We applied \textbf{forward stepwise selection} (AIC-based) to identify
  an optimal set of predictors from a pool of standardized environmental
  variables, including:

  \begin{itemize}
  \tightlist
  \item
    Estimated depth
  \item
    Temperature and salinity at multiple depths
  \item
    Current velocity magnitude and direction
  \item
    Larvae abundance
  \end{itemize}
\end{itemize}

This helped us figure out which factors are most important for
predicting whale presence, and also made it easier to draw inferences
about how each variable might be influencing their behavior.

\subsubsection{Pre-modeling Steps}\label{pre-modeling-steps}

\begin{itemize}
\tightlist
\item
  \textbf{Imputed missing values} in environmental predictors using
  K-nearest neighbors (KNN) to find the nearest (or second-nearest)
  available non-missing value, allowing retention of more samples for
  modeling.
\item
  \textbf{Tested for autocorrelation} in the response using ACF and PACF
  plots.
\item
  \textbf{Split the data temporally}:

  \begin{itemize}
  \tightlist
  \item
    Training set: all years except the most recent
  \item
    Test set: most recent year (to assess temporal generalization)
  \end{itemize}
\end{itemize}

\subsubsection{Model Fitting}\label{model-fitting}

\begin{itemize}
\tightlist
\item
  Both models were fit with:

  \begin{itemize}
  \tightlist
  \item
    A \textbf{spatiotemporal random effect} over season
    (\texttt{spatiotemporal\ =\ "iid"})
  \item
    A \textbf{spatial mesh} created from geographic coordinates
    (\texttt{X}, \texttt{Y})
  \end{itemize}
\item
  Parameters were estimated using \textbf{maximum likelihood} via
  sdmTMB.
\end{itemize}

\subsubsection{Model Evaluation}\label{model-evaluation}

\begin{itemize}
\tightlist
\item
  Performance was assessed by:

  \begin{itemize}
  \tightlist
  \item
    \textbf{Predicted vs.~observed plots} for both training and test
    data
  \item
    \textbf{Root mean squared error (RMSE)}
  \item
    Summary statistics (range and standard deviation of residuals)
  \end{itemize}
\end{itemize}

This modeling framework enabled us to capture non-linear relationships
and spatiotemporal structure in fin and blue whale calling behavior
across oceanographic gradients.

\section{Preliminary Findings}\label{preliminary-findings}

Here are some graphs demonstrating the relationship between predictors
and response with ggplot:

\includegraphics[width=0.8\textwidth,height=\textheight]{ggplotEDA.png}
Using Generalized Additive Model(GAM) to explore the relationship. Here
are some examples below. This is just for a visualized reference,
eventually, we will use variable selection methods (forward step-wise)
to select our final predictors.

\begin{figure}
\centering
\includegraphics[width=0.7\textwidth,height=\textheight]{est.png}
\caption{GAM Fin Whale Response VS. Estimated Depth}
\end{figure}

\begin{figure}
\centering
\includegraphics[width=0.7\textwidth,height=\textheight]{temp_55.png}
\caption{GAM Fin Whale Response VS. Temperature}
\end{figure}

\includegraphics[width=0.7\textwidth,height=\textheight]{magnitude.png}
The GAM plots below reveal clear \textbf{non-linear relationships}
between predictors and whale call activity, supporting the use of
flexible model structures (e.g., Tweedie or delta-gamma families) in
downstream SDM modeling.

\includegraphics[width=0.8\textwidth,height=\textheight]{gif.png}
Through our animations, we were not able to determine anything visually
significant for temperature and salinity across the specified depths.
However, we were able to find a hint about where the whale calls are in
terms of water currents. A majority of the dots which signify how much
effort was recorded, was in the white curves of the graphs. This
represents places in which two different currents are meeting, which
gave us insight into whales' behavior of where to reside and may show
the importance of how the water velocity afffects whale location.

Modeling with the tweedie distribution yielded the following training
and testing plots:

\includegraphics[width=0.5\textwidth,height=\textheight]{train_plot.png}
\includegraphics[width=0.5\textwidth,height=\textheight]{test_plot.png}
Which resulted in a training RMSE of 31.76, a minimum difference of 0
between estimated and predicted values, a maximum difference of 389.22,
and a standard deviation of 31.63. The testing RMSE was 101.19, a
minimum difference of 0, a maximum difference of 728.09 and a standard
deviation of 99.43.

We also fit a Delta Gamma model with yielded the following training and
testing plots:

\includegraphics[width=0.5\textwidth,height=\textheight]{DG_train_plot.png}
\includegraphics[width=0.5\textwidth,height=\textheight]{DG_test_plot.png}

Which resulted in a training RMSE of 14.42, a minimum difference of 0
between estimated and predicted values, a maximum difference of 389.22,
and a standard deviation of 31.63. The testing RMSE was 100.60, a
minimum difference of 0, a maximum difference of 728.09 and a standard
deviation of 99.43.

Interestingly enough, the only difference between these models' metrics
are the RMSE. Both versions have matching maximum and minimum
differences as well as standard deviations. Admittedly, these models
don't do a fantastic job of prediction. With so many of our responses
being that no whales were observed (fin\_resp = 0), our models heavily
overfit toward 0. We even see in our Tweedie model that at times, the
model is predicting negative values. Clearly this cannot be, since
there's no such thing as a negative whale. In the future we hope to find
a way to not overfit so drastically towards 0, but have yet to discover
a solution.

\section{Discussion}\label{discussion}

One of the main challenges we faced was dealing with missing
environmental data. Many variables had non-random gaps, and simply
dropping rows with missing values would have resulted in a major loss of
information. Initially, we considered using spatial interpolation
(kriging), but ultimately, we applied a simpler method: using K-nearest
neighbors (KNN) to find the nearest (or second-nearest) available
non-missing value. It's unclear why missing values appeared during the
extraction process from the CASE-STSE website, but this approach allowed
us to retain more data while maintaining consistency for modeling.

Another major issue came up when we started fitting our models: we ran
into computational singularities, which essentially means the model
couldn't run properly because some variables were too similar to each
other. Specifically, we found that environmental predictors like
salinity at 55 m and 280 m depths were often perfectly correlated. This
kind of high collinearity makes it hard for the model to estimate the
individual effect of each variable. To solve this, we created a
correlation matrix to visualize the relationships between environmental
factors and removed any variables that were highly correlated. We then
used forward stepwise selection to choose a cleaner set of predictors
that wouldn't cause singularities.

We selected the following variables for modeling:

\begin{itemize}
\tightlist
\item
  \texttt{est\_depth\_std}
\item
  \texttt{temperature\_depth\_55\_std}
\item
  \texttt{temperature\_depth\_105\_std}
\item
  \texttt{temperature\_depth\_280\_std}
\item
  \texttt{salinity\_depth\_55\_std}
\item
  \texttt{salinity\_depth\_105\_std}
\item
  \texttt{magnitude\_depth\_55\_std}
\item
  \texttt{magnitude\_depth\_105\_std}
\item
  \texttt{theta\_depth\_55\_std}
\item
  \texttt{theta\_depth\_105\_std}
\item
  \texttt{larvae\_count\_std}
\end{itemize}

\begin{quote}
Note: \texttt{\_std} indicates that the variable has been
\textbf{standardized} (mean-centered and scaled to unit variance).
\end{quote}

Our results from both the tweedie and delta-gamma models showed that the
relationship between environmental factors and fin whale call rates
isn't linear. Factors like depth, temperature, and current strength
seemed to have strong associations with whale calling behavior, which
supports our original idea that environmental conditions help explain
when and where whales are vocalizing.

Finally, we tested our models using the most recent year of data as a
test set. The model did a decent job, but its performance was slightly
worse on the test set compared to the training set, which suggests a bit
of overfitting. In the future, we could address this by adding
regularization or trying out different ways to validate the model.

\section{Future Work}\label{future-work}

In the spring quarter, we plan to build on the foundational work
completed during winter by focusing on completing the three goals that
our sponsor requested of us.

First, we will test whether our tweedie distribution and delta-gamma
distribution models give better performance, and compare models built
with all call types combined versus models that separate call types. We
will continue to refine predictor selection by evaluating variables and
investigating trends. This helps us with inference and understanding how
each predictor affects whale presence.

Second, we will improve our visualizations and draw interpretations,
creating animated maps that overlay whale acoustic presence with
environmental variables, as well as Histograms/scatter plots/x-y plots
to display the relationships of model-predicted whale call density
across time.

Finally, we will begin preparing all project deliverables. This includes
cleaning and documenting the GitHub repository for reproducibility,
designing a poster, and preparing a short presentation for the
showcase--all of which communicates our thoughts on when, where, and why
blue and fin whales are acoustically present in the Southern California
Bight.

\end{document}
